\section{Whole-domain bounds for non-quadratic blocks}\label{app:bounds}

\paragraph{Polynomials.}
Let $f$ be a polynomial with rational coefficients on a rational box
$B=\prod_{i=1}^n[\ell_i,u_i]$. After the affine change of variables
$x_i=\ell_i+(u_i-\ell_i)t_i$ and removal of fixed coordinates, write
\[
f(x(t))=\sum_{\beta\le d}f_\beta t^\beta=\sum_{\theta\le d}b_\theta B^d_\theta(t),
\qquad
B^d_\theta(t)=\prod_i\binom{d_i}{\theta_i}t_i^{\theta_i}(1-t_i)^{d_i-\theta_i},
\]
where $d$ is a componentwise degree bound. The Bernstein coefficients are
\begin{equation}\label{eq:bernstein}
b_\theta=\sum_{\beta\le\theta}f_\beta\prod_i\binom{\theta_i}{\beta_i}\Big/\binom{d_i}{\beta_i}.
\end{equation}

\begin{lemma}[Bernstein enclosure \citep{GarloffSmith2008}]\label{lem:bernstein}
For every $x\in B$, $\min_\theta b_\theta\le f(x)\le\max_\theta b_\theta$.
\end{lemma}

\begin{proof}
On $[0,1]^n$ the basis polynomials are nonnegative and sum to one, so
$f(x(t))$ is a convex combination of the coefficients. Formula
\eqref{eq:bernstein} follows from the univariate identity
$t^j=\sum_{i\ge j}\binom ij\binom dj^{-1}B^d_i(t)$ applied in each coordinate.
\end{proof}

All quantities in~\eqref{eq:bernstein} are rational, so the enclosure is
computed exactly. A weak enclosure is improved by subdivision, and the
enclosures shrink monotonically with the box \citep{GarloffJanssonSmith2003};
exact rational Bernstein subdivision has been formalized in a proof assistant
\citep{MunozNarkawicz2013}.

\begin{proposition}[Partition certificate]\label{prop:partition}
Let $B_1,\dots,B_N$ be closed boxes with $\bigcup_jB_j=B$, and let
$L_j\le\inf_{B_j}f$ for each $j$. Then $\min_jL_j\le\inf_Bf$. Let
$D=\{x\in B:g_r\T x\le\gamma_r,\ r\in\mathcal R\}$ for rational rows, and let $O$
be the set of cells with $\min_{x\in B_j}g_r\T x>\gamma_r$ for some $r$. Then
$\min_{j\notin O}L_j\le\inf_Df$, and if $O$ contains every cell, $D=\emptyset$.
\end{proposition}

The proof is immediate. Its hypothesis is what a checker must verify: the
cells must cover $B$ exactly, starting from the trusted box and not from a
list supplied by the generator. A missing end interval or a child box
generated from a different parent would invalidate the bound. For the same
reason, a stored bound may be reused only for the same expression, direction
and domain.

\paragraph{Univariate elementary expressions.}
For one variable we also admit sums, products and rational powers of $\exp$,
$\log$, $\sin$, $\cos$, $\sinh$, $\cosh$, $\tanh$, $\arctan$ and $\abs{\cdot}$,
with rational constants and $\pi$ and $e$. On every cell of an interval
subdivision, outward-rounded ball arithmetic \citep{Johansson2017arb} gives an
enclosure of the support objective $t\mapsto at+b\T F(t)$. Domain requirements
are part of the expression: $\log$ needs a positive argument, a quotient needs
a denominator interval that excludes zero, and fractional powers use the
nonnegative-base real convention. If an enclosure is not finite or a domain
requirement cannot be proved on a cell, that cell is subdivided further, and
when the depth budget is exhausted the computation is incomplete.

A curvature bound often gives a much better lower bound on a cell than a
natural interval extension. The following bound is the classical remainder of
linear interpolation; with $f''\le M$, the quantity $M(\beta-\alpha)^2/8$ is
also the maximal separation of the $\alpha$BB underestimator of $-f$ with
parameter $M/2$ \citep{AdjimanEtAl1998}.

\begin{lemma}[Chord correction]\label{lem:chord}
Let $f$ be twice continuously differentiable on $[\alpha,\beta]$ with
$f''\le M$. Then
\[
\min_{[\alpha,\beta]}f\ \ge\ \min\{f(\alpha),f(\beta)\}-\max\{0,M\}\,(\beta-\alpha)^2/8.
\]
\end{lemma}

\begin{proof}
The function $f(t)-Mt^2/2$ is concave, so it lies above its chord. Rearranging
gives $f(t)\ge\ell_f(t)-\tfrac M2(t-\alpha)(\beta-t)$, where $\ell_f$ is the chord
of $f$. If $M\ge0$, use $(t-\alpha)(\beta-t)\le(\beta-\alpha)^2/4$ and
$\ell_f(t)\ge\min\{f(\alpha),f(\beta)\}$. If $M<0$, the correction term is
nonnegative and $f\ge\ell_f$.
\end{proof}

Verified enclosures of $f(\alpha)$, $f(\beta)$ and of $f''$ on the cell are
enough, provided that $f$ is twice continuously differentiable on the cell.
On a cell that contains a nondifferentiable point of the expression, such as
a zero of the argument of $\abs{\cdot}$, only the natural enclosure is used.
The implementation takes the larger of the two bounds on each cell.
